\documentclass[]{book}
\usepackage{xcolor}
\usepackage{listings}
\definecolor{ghblack}{HTML}{24292E}
\definecolor{ghblue}{HTML}{0366D6}
\definecolor{ghorange}{HTML}{E36209}
\definecolor{ghpurple}{HTML}{6F42C1}
\definecolor{ghgreen}{HTML}{22863A}
\definecolor{ghgray}{HTML}{6A737D}
\definecolor{ghbg}{HTML}{F6F8FA}
\lstdefinestyle{githubpy}{
    language=Python,
    backgroundcolor=\color{ghbg},
    basicstyle=\ttfamily\footnotesize\color{ghblack},
    commentstyle=\color{ghgray},
    keywordstyle=\color{ghpurple},
    stringstyle=\color{ghblue},
    numberstyle=\tiny\color{ghgray},
    breakatwhitespace=false,
    breaklines=true,
    captionpos=b,
    keepspaces=true,
    numbers=left,
    numbersep=5pt,
    showspaces=false,
    showstringspaces=false,
    showtabs=false,
    tabsize=4,
    frame=single,
    framesep=5pt,
    framerule=0pt,
    rulecolor=\color{ghgray},
    identifierstyle=\color{ghblack},
    emphstyle=\color{ghorange},
    emph={self},
    morekeywords={as, assert, break, class, continue, def, del, elif, else, except, finally, for, from, global, if, import, in, is, lambda, nonlocal, not, or, pass, raise, return, try, while, with, yield},
}

%These tell TeX which packages to use.
\usepackage{array,epsfig}
\usepackage{amsmath}
\usepackage{amsfonts}
\usepackage{amssymb}
\usepackage{amsxtra}
\usepackage{amsthm}
\usepackage{mathrsfs}
\usepackage{color}
\usepackage{graphicx}
\graphicspath{ {./images/} }
%Here I define some theorem styles and shortcut commands for symbols I use often
\theoremstyle{definition}
\newtheorem{defn}{Definition}
\newtheorem{thm}{Theorem}
\newtheorem{cor}{Corollary}
\newtheorem*{rmk}{Remark}
\newtheorem{lem}{Lemma}
\newtheorem*{joke}{Joke}
\newtheorem{ex}{Example}
\newtheorem*{soln}{Solution}
\newtheorem{prop}{Proposition}

\newcommand{\lra}{\longrightarrow}
\newcommand{\ra}{\rightarrow}
\newcommand{\surj}{\twoheadrightarrow}
\newcommand{\graph}{\mathrm{graph}}
\newcommand{\bb}[1]{\mathbb{#1}}
\newcommand{\Z}{\bb{Z}}
\newcommand{\Q}{\bb{Q}}
\newcommand{\R}{\bb{R}}
\newcommand{\C}{\bb{C}}
\newcommand{\N}{\bb{N}}
\newcommand{\M}{\mathbf{M}}
\newcommand{\m}{\mathbf{m}}
\newcommand{\MM}{\mathscr{M}}
\newcommand{\HH}{\mathscr{H}}
\newcommand{\Om}{\Omega}
\newcommand{\Ho}{\in\HH(\Om)}
\newcommand{\bd}{\partial}
\newcommand{\del}{\partial}
\newcommand{\bardel}{\overline\partial}
\newcommand{\textdf}[1]{\textbf{\textsf{#1}}\index{#1}}
\newcommand{\img}{\mathrm{img}}
\newcommand{\ip}[2]{\left\langle{#1},{#2}\right\rangle}
\newcommand{\inter}[1]{\mathrm{int}{#1}}
\newcommand{\exter}[1]{\mathrm{ext}{#1}}
\newcommand{\cl}[1]{\mathrm{cl}{#1}}
\newcommand{\ds}{\displaystyle}
\newcommand{\vol}{\mathrm{vol}}
\newcommand{\cnt}{\mathrm{ct}}
\newcommand{\osc}{\mathrm{osc}}
\newcommand{\LL}{\mathbf{L}}
\newcommand{\UU}{\mathbf{U}}
\newcommand{\support}{\mathrm{support}}
\newcommand{\AND}{\;\wedge\;}
\newcommand{\OR}{\;\vee\;}
\newcommand{\Oset}{\varnothing}
\newcommand{\st}{\ni}
\newcommand{\wh}{\widehat}

%Pagination stuff.
\setlength{\topmargin}{-.3 in}
\setlength{\oddsidemargin}{0in}
\setlength{\evensidemargin}{0in}
\setlength{\textheight}{9.in}
\setlength{\textwidth}{6.5in}
\pagestyle{empty}



\begin{document}
\subsection*{Low-rank compression for Q}
\paragraph{}
Given $h_t \in \mathbb{R}^{B \times L \times 5120}$, where $B$ is the batch size and $L$ is the sequence length. In DeepSeek-V2, the Q vector also adopts low-rank compression. First, the input vector is projected to a 1536-dimensional low-dimensional space:
\begin{equation}
    c_t^Q = W^{DQ} h_t \in \mathbb{R}^{B \times L \times 1536}
\end{equation}
\paragraph{}
where, $W^{DQ} \in \mathbb{R}^{d_{c} \times d}$. $d$ is the hidden dimension ($5120$ in this case) and $d_{c}$ is the query compression dimension ($1536$ in this case)
\paragraph{}
Then, it is projected (decompression) to the $\mathbb{R}^{H \times 128}$ multi-head vector space (where $H=128$ is the number of heads), obtaining the first part of the Q vector:
\begin{equation}
    q_t^C = W^{UQ} c_t^Q \in \mathbb{R}^{B \times L \times H \times 128}
\end{equation}
\paragraph{}
Next, it is projected to $\mathbb{R}^{H \times 64}$ and uses RoPE to embed position information, obtaining the second part of the Q vector:
\begin{equation}
    q_t^R = \mathrm{RoPE}(W^{KR} h_t) \in \mathbb{R}^{B \times L \times H \times 64} 
\end{equation}
\paragraph{}
The two parts are concatenated to obtain the final Q vector:
\begin{equation}
    q_t = [q_t^C, q_t^R] \in \mathbb{R}^{B \times L \times H \times 192}
\end{equation}
\begin{lstlisting}[style=githubpy]
q = self.q_b_proj(self.q_a_layernorm(self.q_a_proj(hidden_states)))
q = q.view(bsz, q_len, self.num_heads, self.q_head_dim).transpose(1, 2)
\end{lstlisting}
\subsection*{Low-rank joint compression for K and V}
\paragraph{}
When computing the KV vector, the input vector first needs to be projected to a 512-dimensional joint compressed representation:
\begin{equation}
    c_t^{KV} = W^{DKV} h_t \in \mathbb{R}^{B \times L \times 512}
\end{equation}
\paragraph{}
Similar to the computation process of the Q vector, the first part of the K vector is obtained by projecting $c_t^{KV}$ through decompression to the $\mathbb{R}^{H \times 128}$ multi-head vector space:
\begin{equation}
    k_t^C = W^{UK} c_t^{KV} \in \mathbb{R}^{B \times L \times H \times 128}
\end{equation}
\paragraph{}
The second part of K is obtained by projecting the input vector to a 64-dimensional vector space and applying RoPE to embed position information:
\begin{equation}
    k_t^R = \mathrm{RoPE}(W^{KR} h_t) \in \mathbb{R}^{B \times L \times 64}
\end{equation}
\paragraph{}
Unlike Q, the complete K is obtained by broadcasting the second part of K to each head and concatenating it with the first part:
\begin{equation}
     k_t = \begin{bmatrix}
    k_{t,1}^C & k_t^R \\ 
    k_{t,2}^C & k_t^R \\
    \vdots & \vdots \\
    \end{bmatrix} \in \mathbb{R}^{B \times L \times H \times 192}
\end{equation}
\paragraph{}
That is, the RoPE part of each head is exactly the same.
\paragraph{}
The computation of the V vector is relatively simple, directly decompressing $c_t^{KV}$ to $\mathbb{R}^{H \times 128}$:
\begin{equation}
    v_t = W^{UV} c_t^{KV} \in \mathbb{R}^{B \times L \times H \times 128} 
\end{equation}
\paragraph{}
The attention computation process is no different from traditional MHA. First, the attention score is computed:
\begin{equation}
a = \mathrm{softmax}\left(\frac{q_t^\top k_t + \mathrm{Mask}}{\sqrt{192}}\right) = 
\mathrm{softmax}\left(\frac{{q_t^C}^\top k_t^C + {q_t^R}^\top k_t^R + \mathrm{Mask}}{\sqrt{128 + 64}} \right)
\in \mathbb{R}^{B \times L \times H \times L}
\end{equation}
\paragraph{}
The weighted sum of V is computed, and all heads are flattened to obtain the Attention output:
\begin{equation}
    o = a \cdot v_t \in \mathbb{R}^{B \times L \times H \times 128} \cong \mathbb{R}^{B \times L \times 16384}
\end{equation}
\paragraph{}
After projection with another matrix, the final output of MLA is obtained:
\begin{equation}
    u = W^O o \in \mathbb{R}^{B \times L \times 5120}
\end{equation}
\paragraph{}
\begin{lstlisting}[style=githubpy]
def forward(...):
    bsz, q_len, _ = hidden_states.size()
    
    # Compute Q: first reduce dimension then increase dimension
    # Compared to directly using a matrix of size [5120, 24576], 
    # the low-rank decomposition [5120, 1536] * [1536, 24576] greatly reduces storage space and computation
    q = self.q_b_proj(self.q_a_layernorm(self.q_a_proj(hidden_states)))
    q = q.view(bsz, q_len, self.num_heads, self.q_head_dim).transpose(1, 2)
    # Split rope and non-rope parts
    q_nope, q_pe = torch.split(
        q, [self.qk_nope_head_dim, self.qk_rope_head_dim], dim=-1
    )
    
    # Compute KV
    # An optimized MLA KVCache implementation only needs to cache this compressed_kv, but it is actually expanded later
    compressed_kv = self.kv_a_proj_with_mqa(hidden_states)
    # Here compressed_kv corresponds to c_t^{KV} in the formula
    compressed_kv, k_pe = torch.split(
        compressed_kv, [self.kv_lora_rank, self.qk_rope_head_dim], dim=-1
    )
    k_pe = k_pe.view(bsz, q_len, 1, self.qk_rope_head_dim).transpose(1, 2)
    # Expand MLA to the standard MHA form
    kv = (
        self.kv_b_proj(self.kv_a_layernorm(compressed_kv))
        .view(bsz, q_len, self.num_heads, self.qk_nope_head_dim + self.v_head_dim)
        .transpose(1, 2)
    )
    # Since kv_b_proj packs W^{UK} and W^{UV} together, separate them
    k_nope, value_states = torch.split(
        kv, [self.qk_nope_head_dim, self.v_head_dim], dim=-1
    )
    ...
    # Add rope to the parts that need rope
    q_pe, k_pe = apply_rotary_pos_emb(q_pe, k_pe, cos, sin, position_ids)
    
    # Update and concatenate historical KVCache
    # It can be seen that the expanded MHA KVCache is stored here
    # where q_head_dim equals qk_nope_head_dim + qk_rope_head_dim
    query_states = k_pe.new_empty(bsz, self.num_heads, q_len, self.q_head_dim)
    query_states[:, :, :, : self.qk_nope_head_dim] = q_nope
    query_states[:, :, :, self.qk_nope_head_dim :] = q_pe
    key_states = k_pe.new_empty(bsz, self.num_heads, q_len, self.q_head_dim)
    key_states[:, :, :, : self.qk_nope_head_dim] = k_nope
    key_states[:, :, :, self.qk_nope_head_dim :] = k_pe
    if past_key_value is not None:
        cache_kwargs = {"sin": sin, "cos": cos}  # Specific to RoPE models
        key_states, value_states = past_key_value.update(
            key_states, value_states, self.layer_idx, cache_kwargs
        )

    # The subsequent code is the standard MHA code, no need to elaborate
    ...
\end{lstlisting}
\end{document}



